\documentclass[conference]{IEEEtran}
\IEEEoverridecommandlockouts

\usepackage[utf8]{inputenc}
\usepackage[T1]{fontenc}
\usepackage[spanish,mexico]{babel}
\usepackage{cite}
\usepackage{amsmath,amssymb}
\usepackage{graphicx}
\usepackage{booktabs}
\usepackage{url}
\usepackage{hyperref}

\begin{document}

\title{Predicción de precios y clasificación de sistemas constructivos\\
en el mercado inmobiliario de Tokio:\\
un estudio sobre 344.000 transacciones}

\author{
\IEEEauthorblockN{Nombre del Autor}
\IEEEauthorblockA{\textit{Licenciatura en Ciencia de Datos} \\
\textit{Pontificia Universidad Católica Argentina}\\
Rosario, Argentina \\
correo@ejemplo.edu.ar}
}

\maketitle

\begin{abstract}
El mercado inmobiliario de Tokio presenta una heterogeneidad extrema: conviven operaciones
residenciales de pocos millones de yenes con transacciones corporativas que superan los
\textyen170.000 millones. Este trabajo aborda dos tareas complementarias sobre el registro
público de transacciones del Ministerio de Tierra, Infraestructura, Transporte y Turismo
(MLIT) de Japón: la estimación del valor de transacción y la clasificación del sistema
constructivo del edificio. Se procesaron 343.995 operaciones de la Prefectura de Tokio
(2005--2025) resolviendo tres problemas de calidad no triviales: un artefacto de codificación
en las superficies, valores textuales en la distancia a estación y un objetivo de
clasificación multi-etiqueta. Para la regresión se compararon tres modelos sobre el logaritmo
del precio, obteniendo XGBoost el mejor desempeño ($R^2_{\log}=0{,}848$). Para la
clasificación, con un desbalance de 65,6:1, XGBoost alcanzó un F1 macro de 0,566 y un ROC-AUC
one-vs-rest de 0,948. Se identifica una relación contraintuitiva: la distancia a la estación
correlaciona débilmente con el precio ($r=-0{,}07$) pero fuertemente con los índices de
edificabilidad, sugiriendo que su efecto opera por vía de la zonificación. Ambos modelos se
desplegaron en una aplicación web que comunica explícitamente la incertidumbre de sus
estimaciones.
\end{abstract}

\begin{IEEEkeywords}
aprendizaje automático, mercado inmobiliario, regresión, clasificación multiclase, desbalance
de clases, XGBoost
\end{IEEEkeywords}

\section{Introducción}

La valuación inmobiliaria automatizada (\textit{Automated Valuation Model}, AVM) constituye una
de las aplicaciones más consolidadas del aprendizaje automático sobre datos tabulares. Su
utilidad abarca desde la tasación asistida hasta el análisis de política urbana, y su
viabilidad depende críticamente de la disponibilidad de registros de transacciones
suficientemente extensos y confiables.

El caso japonés resulta particularmente interesante. El MLIT publica de forma abierta los
precios de transacción reportados por las partes intervinientes, con un nivel de detalle que
incluye atributos del lote, del edificio y del entorno urbano \cite{mlit}. Sin embargo, la
Prefectura de Tokio presenta una dificultad estructural: su mercado es extraordinariamente
heterogéneo. Los distritos centrales concentran edificios corporativos cuyo valor supera en
varios órdenes de magnitud al de las viviendas unifamiliares de los distritos periféricos.

Trabajos previos sobre predicción de precios inmobiliarios han mostrado que los métodos de
ensamble basados en árboles superan consistentemente a los modelos lineales en datos
tabulares heterogéneos \cite{chen2016,breiman2001}. La transformación logarítmica del objetivo
es práctica habitual ante distribuciones de precios con asimetría severa \cite{hastie2009}.
Por su parte, la clasificación bajo desbalance de clases exige métricas específicas: la
exactitud resulta engañosa cuando una clase domina la muestra, y se recomienda el uso de
F1 macro o coeficiente kappa \cite{he2009}.

Este trabajo persigue tres objetivos:

\begin{enumerate}
\item Caracterizar y resolver los problemas de calidad del registro del MLIT para la
      Prefectura de Tokio, documentando cada decisión metodológica.
\item Comparar tres algoritmos de regresión para estimar el valor de transacción, y tres
      algoritmos de clasificación para predecir el sistema constructivo del edificio.
\item Desplegar los modelos resultantes en una aplicación web que comunique de manera
      responsable la incertidumbre de sus predicciones.
\end{enumerate}

\section{Metodología}

\subsection{Datos}

El conjunto de datos comprende \textbf{343.995 transacciones} de la Prefectura de Tokio
registradas entre 2005 y 2025, descritas por 27 variables originales. Cada fila corresponde a
una operación de compraventa reportada al MLIT.

Una decisión de diseño temprana fue trabajar sobre la prefectura completa en lugar de acotarse
a un único distrito. Esto habilita el uso del \textit{ward} (distrito administrativo) como
variable predictora, incorporando 59 categorías geográficas que, como muestra la
Sección~\ref{sec:eda}, resultan determinantes para el precio.

\subsection{Problemas de calidad y su tratamiento}

El preprocesamiento debió resolver cinco problemas no triviales.

\textbf{Artefacto de codificación.} Las columnas de superficie llegan en codificación
Windows-1252, donde el carácter japonés correspondiente a metro cuadrado se corrompe en una
secuencia espuria de dos caracteres ASCII. Los valores del extremo superior aparecen como
\texttt{2,000 [artefacto] or greater.}. Una conversión numérica directa los habría convertido en
nulos, eliminando precisamente las operaciones de mayor superficie. Se extrajo el valor
mediante expresiones regulares y se marcó la censura con una variable indicadora.

\textbf{Distancia a estación no numérica.} La variable combina minutos numéricos con rangos
textuales (\texttt{30-60minutes}, \texttt{1H-1H30}, \texttt{2H-}). Se implementó un
analizador que convierte cada formato a minutos, tomando el punto medio en los intervalos
cerrados y el extremo inferior en los abiertos.

\textbf{Nulos estructurales.} Las operaciones de tipo \textit{Land Only} carecen de edificio
por definición: superficie construida, año de construcción y estructura están vacíos por
naturaleza del registro, no por ausencia de dato. Imputar estadísticamente esas columnas
habría fabricado edificios inexistentes. Se optó por un indicador binario
\texttt{sin\_edificio} y relleno neutro.

\textbf{Inconsistencias temporales.} Se detectaron 7.066 registros con antigüedad negativa
--- edificios cuyo año de construcción es posterior al de la operación ---. Se marcaron con
una variable indicadora y se llevaron a cero.

\textbf{Objetivo multi-etiqueta.} La variable \texttt{Building : Structure} no presenta cuatro
categorías limpias sino aproximadamente 25, incluyendo combinaciones (\texttt{RC, W}) y
clases marginales (\texttt{LS}, \texttt{B}). Se derivó \texttt{estructura\_principal}
conservando el primer sistema constructivo listado y agrupando las categorías marginales bajo
\texttt{Otros}.

\subsection{Ingeniería de características}

Se derivaron variables temporales (año y trimestre de operación, antigüedad del edificio), de
densidad (\texttt{ratio\_construido} como cociente entre superficie construida y del terreno) y
seis indicadores binarios de uso a partir de la columna multi-valor \texttt{Use}.

El tratamiento de la cardinalidad geográfica mereció atención específica. El \textit{ward}
(59 categorías) se codificó mediante \textit{one-hot}, por ser interpretable y de cardinalidad
manejable. En cambio \texttt{District} (1.449 categorías) y la estación más cercana (660) se
redujeron a las 100 más frecuentes agrupando el resto bajo \texttt{Otros}, para luego aplicar
codificación por frecuencia. Cabe señalar que un criterio inicial basado en umbral porcentual
(0,5\%) resultó inadecuado: con 344.000 filas exigía más de 1.700 casos por categoría,
relegando el 98\% de los registros a \texttt{Otros}. El criterio top-$N$ controla el número de
categorías con independencia del tamaño de la muestra.

\subsection{Diseño experimental}

Ambas tareas emplean partición \textit{hold-out} 80/20 con semilla fija y validación cruzada de
5 particiones sobre el conjunto de entrenamiento. Todo el preprocesamiento se encapsuló en
objetos \texttt{Pipeline} de scikit-learn \cite{pedregosa2011}, de modo que la imputación y el
escalado se ajustan exclusivamente sobre los datos de entrenamiento en cada partición,
evitando fuga de información.

Para la regresión, el objetivo se modela como $\log(1+y)$ y las métricas se reportan en yenes
aplicando la transformación inversa. Se excluyeron del conjunto de predictores las variables
derivadas del precio (\texttt{precio\_suelo\_m2}, \texttt{precio\_por\_m2}) por constituir fuga
directa del objetivo.

Para la clasificación se aplicó partición estratificada y \texttt{class\_weight="balanced"} en
los modelos que lo admiten. La métrica de referencia es el F1 macro.

\section{Resultados}
\label{sec:eda}

\subsection{Análisis exploratorio}

La distribución del precio presenta una asimetría extrema (coeficiente de asimetría de 150,4),
con mediana de \textyen46 millones, media de \textyen96,5 millones y máximo de
\textyen170.000 millones. La transformación logarítmica reduce la asimetría a 0,38,
justificando su adopción como escala de modelado.

El hallazgo más relevante del análisis exploratorio proviene de la comparación entre distritos.
El precio mediano por metro cuadrado varía en un factor cercano a cuatro: Chiyoda alcanza
\textyen2,4 millones/m\textsuperscript{2} mientras que Sumida se sitúa en torno a
\textyen0,63 millones/m\textsuperscript{2}. Más aún, el ordenamiento por volumen de operaciones
resulta prácticamente inverso: Setagaya, Nerima y Adachi lideran en cantidad de transacciones
sin figurar entre los distritos más caros. Esta disociación entre precio y volumen justifica
empíricamente la incorporación del \textit{ward} como predictor.

Un segundo hallazgo resulta contraintuitivo. La distancia a la estación ferroviaria muestra una
correlación muy débil con el precio ($r = -0{,}07$), contra la expectativa de que la
accesibilidad fuera un determinante principal. Sin embargo, correlaciona apreciablemente con
los índices de edificabilidad: $-0{,}36$ con el \textit{building coverage ratio} y $-0{,}33$
con el \textit{floor area ratio}. La interpretación plausible es que la proximidad a estaciones
opera principalmente a través de la zonificación urbana --- las áreas mejor conectadas admiten
mayor densidad constructiva --- y que su efecto sobre el precio es por tanto indirecto.

Se detectó además colinealidad elevada entre \textit{building coverage ratio} y \textit{floor
area ratio} ($r = 0{,}81$), factor a considerar en la interpretación del modelo lineal.

\subsection{Regresión sobre el precio}

La Tabla~\ref{tab:regresion} resume el desempeño de los tres modelos.

\begin{table}[htbp]
\caption{Desempeño de los modelos de regresión (conjunto de prueba)}
\label{tab:regresion}
\centering
\begin{tabular}{lrrrr}
\toprule
\textbf{Modelo} & $\mathbf{R^2_{\log}}$ & $\mathbf{R^2_{\text{\textyen}}}$ & \textbf{MAE (\textyen)} & \textbf{MAPE} \\
\midrule
Regresión Lineal & 0,724 & 0,100 & 47.182.148 & 164,3\% \\
Random Forest    & 0,705 & 0,077 & 43.225.680 & 195,9\% \\
\textbf{XGBoost} & \textbf{0,848} & \textbf{0,267} & \textbf{29.490.866} & \textbf{140,1\%} \\
\bottomrule
\end{tabular}
\end{table}

XGBoost domina en todas las métricas evaluadas. La superioridad frente a la regresión lineal
(0,848 contra 0,724 en $R^2_{\log}$) confirma la presencia de relaciones no lineales e
interacciones entre ubicación, superficie y antigüedad que el modelo lineal no captura.

Merece discusión la brecha entre el $R^2$ calculado en escala logarítmica (0,848) y el
calculado en yenes (0,267). Ambos describen el mismo modelo: la diferencia obedece a que la
transformación inversa $\exp(x)-1$ amplifica exponencialmente los errores en la cola superior
de la distribución. Un error de 0,37 en escala logarítmica --- el desvío observado de los
residuos --- se traduce en miles de millones de yenes cuando se aplica sobre una operación de
\textyen170.000 millones. El $R^2$ logarítmico constituye la comparación justa entre modelos,
puesto que es la escala en que estos optimizan; el $R^2$ en yenes informa sobre la utilidad
práctica de las estimaciones en valores absolutos.

El análisis de importancia de variables sitúa a \texttt{area\_construida\_m2} entre los
predictores principales, junto con múltiples indicadores de \textit{ward}, lo que respalda
la decisión de incorporar la dimensión geográfica.

\subsection{Clasificación del sistema constructivo}

La distribución del objetivo evidencia un desbalance severo: \texttt{W} (madera) representa el
80,8\% de los casos, seguida de \texttt{S} (7,5\%), \texttt{RC} (7,45\%), \texttt{Otros}
(3,0\%) y \texttt{SRC} (1,2\%), con una relación de 65,6:1 entre la clase mayoritaria y la
minoritaria.

\begin{table}[htbp]
\caption{Desempeño de los modelos de clasificación (conjunto de prueba)}
\label{tab:clasificacion}
\centering
\begin{tabular}{lrrrr}
\toprule
\textbf{Modelo} & \textbf{F1 macro} & \textbf{Exactitud} & \textbf{Kappa} & \textbf{ROC-AUC} \\
\midrule
Reg. Logística   & 0,493 & 0,725 & 0,409 & 0,903 \\
Random Forest    & 0,554 & 0,772 & 0,485 & 0,932 \\
\textbf{XGBoost} & \textbf{0,566} & \textbf{0,892} & \textbf{0,635} & \textbf{0,948} \\
\bottomrule
\end{tabular}
\end{table}

La Tabla~\ref{tab:clasificacion} ilustra por qué la exactitud resulta inadecuada como métrica
de referencia. XGBoost alcanza un 89,2\% de exactitud, pero un clasificador trivial que
predijera siempre \texttt{W} obtendría aproximadamente 81\% sin detectar ningún otro sistema
constructivo. El F1 macro de 0,566 ofrece una lectura más honesta del desempeño.

El desempeño por clase confirma esta interpretación: el \textit{recall} sobre \texttt{W}
alcanza 0,987, mientras que sobre \texttt{Otros} se reduce a 0,048 y sobre \texttt{SRC} a
0,454. Como se anticipaba, la matriz de confusión revela que \texttt{RC} y \texttt{SRC}
concentran la mayor confusión mutua, resultado esperable dado que ambos sistemas comparten
características constructivas y perfiles de uso similares.

\subsection{Análisis no supervisado}

Se aplicó agrupamiento por $k$-medias sobre las variables preprocesadas, con $k$ igual al
número de clases. El índice de Rand ajustado entre las agrupaciones obtenidas y las clases
reales alcanzó 0,211, valor que indica una coincidencia parcial pero claramente insuficiente.

Este resultado admite una lectura sustantiva: la estructura natural de los datos responde
principalmente a la ubicación y el tamaño de los inmuebles, no a su sistema constructivo. El
material de un edificio depende de factores --- normativa antisísmica vigente al momento de la
construcción, época, decisión del desarrollador --- que las variables disponibles capturan solo
de manera parcial. Esta limitación explica también el techo observado en el desempeño de los
clasificadores supervisados.

\section{Despliegue}

Los modelos se integraron en una aplicación web desarrollada con Streamlit, organizada en dos
pestañas correspondientes a cada tarea. Los artefactos se serializan como objetos
\texttt{Pipeline} completos, de modo que la aplicación no replica lógica de preprocesamiento:
recibe los datos del formulario en su forma cruda y delega las transformaciones en el pipeline.

Un aspecto central del diseño fue la comunicación de la incertidumbre. Dado el MAPE del 140\%
documentado en la Tabla~\ref{tab:regresion}, presentar una estimación puntual induciría una
falsa sensación de precisión. La aplicación muestra en cambio un intervalo construido a partir
del desvío de los residuos en escala logarítmica, necesariamente asimétrico al expresarse en
yenes, y emite una advertencia explícita cuando la estimación supera el umbral a partir del
cual el error esperado se incrementa sustancialmente.

De modo análogo, la pestaña de clasificación presenta la distribución completa de
probabilidades por clase y señala cuando la predicción corresponde a una categoría minoritaria
sobre la cual el modelo exhibe bajo \textit{recall}.

\section{Discusión y conclusiones}

Este trabajo confirma que los métodos de ensamble basados en árboles superan a los modelos
lineales en datos inmobiliarios tabulares, en línea con la literatura previa. La mejora de
XGBoost sobre la regresión lineal (0,848 contra 0,724 en $R^2_{\log}$) es consistente y se
sostiene en todas las métricas evaluadas.

La etapa más desafiante no fue el modelado sino la preparación de los datos. Tres de los cinco
problemas de calidad identificados --- el artefacto de codificación, la distancia textual y el
objetivo multi-etiqueta --- habrían pasado inadvertidos bajo un procesamiento superficial, y
cada uno tenía potencial para invalidar los resultados. El caso del objetivo multi-etiqueta es
ilustrativo: sin normalización, el clasificador habría entrenado sobre aproximadamente 25
categorías, varias con un puñado de ejemplos.

Respecto de la utilidad práctica, corresponde ser cauto. El modelo de regresión resulta
adecuado para estimar el orden de magnitud de una operación típica, pero su MAPE del 140\% lo
inhabilita como herramienta de tasación. El clasificador ofrece un desempeño razonable sobre
las clases frecuentes y limitado sobre las minoritarias.

Deben señalarse asimismo las limitaciones éticas y de sesgo. Los datos son autoreportados por
las partes intervinientes, lo que admite subreporte sistemático. La cobertura se restringe a
Tokio, de modo que el modelo no generaliza a otras prefecturas. Y el registro omite atributos
determinantes del valor --- calidad constructiva, estado de conservación, condiciones
particulares de cada negociación --- que ningún algoritmo puede inferir de las variables
disponibles. Un sistema de valuación desplegado sin estas advertencias podría inducir
decisiones económicas mal fundadas.

Como líneas de trabajo futuro se identifican: la incorporación de variables geoespaciales
explícitas (coordenadas, distancia al centro), la evaluación de técnicas de sobremuestreo para
las clases minoritarias, y la construcción de intervalos de predicción mediante regresión
cuantílica en lugar de la aproximación por desvío de residuos empleada aquí.

\begin{thebibliography}{00}

\bibitem{mlit}
Ministry of Land, Infrastructure, Transport and Tourism of Japan, ``Real Estate Transaction
Price Information,'' 2025. [En línea]. Disponible:
\url{https://www.land.mlit.go.jp/webland_english/servlet/MainServlet}

\bibitem{chen2016}
T. Chen y C. Guestrin, ``XGBoost: A scalable tree boosting system,'' en \textit{Proc. 22nd ACM
SIGKDD Int. Conf. Knowledge Discovery and Data Mining}, San Francisco, CA, EE.UU., 2016,
pp. 785--794.

\bibitem{breiman2001}
L. Breiman, ``Random forests,'' \textit{Machine Learning}, vol. 45, n.\textsuperscript{o} 1,
pp. 5--32, 2001.

\bibitem{hastie2009}
T. Hastie, R. Tibshirani y J. Friedman, \textit{The Elements of Statistical Learning: Data
Mining, Inference, and Prediction}, 2.\textsuperscript{a} ed. Nueva York, NY, EE.UU.:
Springer, 2009.

\bibitem{he2009}
H. He y E. A. Garcia, ``Learning from imbalanced data,'' \textit{IEEE Trans. Knowledge and
Data Engineering}, vol. 21, n.\textsuperscript{o} 9, pp. 1263--1284, 2009.

\bibitem{pedregosa2011}
F. Pedregosa \textit{et al.}, ``Scikit-learn: Machine learning in Python,'' \textit{Journal of
Machine Learning Research}, vol. 12, pp. 2825--2830, 2011.

\end{thebibliography}

\end{document}
